\documentclass[conference]{IEEEtran}
\usepackage{float}
\IEEEoverridecommandlockouts

\usepackage{amsmath,amssymb,amsfonts}
\usepackage{graphicx}
\usepackage{booktabs}
\usepackage{xcolor}
\usepackage{url}
\usepackage{cite}
\usepackage{hyperref}
\usepackage{placeins}
\usepackage{dblfloatfix}
\usepackage{stfloats}
\hypersetup{colorlinks=true,linkcolor=blue,citecolor=blue,urlcolor=blue}

\renewcommand{\topfraction}{0.95}
\renewcommand{\bottomfraction}{0.95}
\renewcommand{\textfraction}{0.05}
\renewcommand{\floatpagefraction}{0.80}
\renewcommand{\dbltopfraction}{0.95}
\renewcommand{\dblfloatpagefraction}{0.80}
\setcounter{topnumber}{4}
\setcounter{bottomnumber}{3}
\setcounter{totalnumber}{6}
\setcounter{dbltopnumber}{4}

\begin{document}

\title{Hybrid CNN-LSTM Model for Short-Term Solar Power Forecasting with Anomaly-Aware Preprocessing}

\author{
\IEEEauthorblockN{Tushar Basak}
\IEEEauthorblockA{\textit{Department of Electronics \& Communication Engineering} \\
\textit{Faculty of Technology, University of Delhi}\\
Delhi 110007, India \\
tusharbasak555@ece.du.ac.in}
\and
\IEEEauthorblockN{Varun Saxena}
\IEEEauthorblockA{\textit{Department of Electronics \& Communication Engineering} \\
\textit{Faculty of Technology, University of Delhi}\\
Delhi 110007, India \\
varunsaxena@fot.du.ac.in}
}

\maketitle

\begin{abstract}
Accurate ultra-short-term photovoltaic (PV) power forecasting depends not only on forecasting architecture but also on the quality of plant telemetry. This paper investigates an anomaly-aware preprocessing pipeline for 30-minute-ahead PV power forecasting. Using a one-year, 5-minute-resolution NREL dataset from a 70~MW utility-scale PV site, we synthetically inject three classes of telemetry corruption -- daytime dropouts, spike anomalies, and short-duration offset faults -- across $\sim$2\% of observations. An Isolation Forest identifies anomalous observations, which are reconstructed through local linear interpolation. To isolate the contribution of preprocessing from that of detection quality, we evaluate four parallel data conditions: a never-corrupted \emph{pristine} reference (the performance ceiling), \emph{raw corrupted} telemetry, telemetry cleaned by the unsupervised \emph{Isolation Forest} detector (the proposed pipeline), and telemetry cleaned using \emph{oracle} ground-truth corruption labels. The detector achieves 72.21\% recall and 57.48\% precision. All deep learning results are reported as mean $\pm$ standard deviation across 5 random seeds, with paired $t$-tests, a Holm--Bonferroni correction for multiple comparisons, and a Two One-Sided Tests (TOST) equivalence analysis used to validate every comparative claim. The proposed CNN--LSTM achieves RMSE $= 3.57 \pm 0.04$~MW, MAE $= 1.64 \pm 0.08$~MW, and $R^2 = 0.949 \pm 0.001$. Paired tests against the pristine ceiling ($3.58 \pm 0.06$~MW) and the oracle-cleaned condition ($3.58 \pm 0.05$~MW) show no significant difference (Holm-adjusted $p = 1.000$ for both), and a TOST equivalence test confirms the proposed model is statistically indistinguishable from the pristine ceiling within a $\pm 0.10$~MW practical-equivalence margin ($p \approx 0.05$). Against the raw-corrupted ablation ($3.89 \pm 0.11$~MW), cleaning yields an 8.2\% mean RMSE reduction. The proposed model achieves a lower RMSE than a plain LSTM baseline in every one of 5 paired seeds ($3.77 \pm 0.11$~MW; raw paired $t=-3.89$, $p=0.018$, Holm adjusted $p=0.053$). Both deep models outperform persistence (4.93~MW) and a diurnal-lag autoregressive baseline (4.83~MW) by a wide margin.
\end{abstract}

\begin{IEEEkeywords}
Anomaly detection, anomaly-aware preprocessing, CNN--LSTM, Isolation Forest, photovoltaic power forecasting, sensor fault detection, short-term forecasting, time-series imputation, equivalence testing.
\end{IEEEkeywords}

% =================================================================
\section*{Nomenclature}
\addcontentsline{toc}{section}{Nomenclature}
\noindent
\begin{IEEEdescription}[\IEEEsetlabelwidth{SCADA}]
\item[IF] Isolation Forest
\item[SCADA] Supervisory Control and Data Acquisition
\item[TOST] Two One-Sided Tests (equivalence testing)
\item[UPV] Utility-Scale Photovoltaic
\item[$H$] Forecast horizon (time steps)
\item[$W$] Input window length (time steps)
\item[$F$] Number of input features per time step
\item[$P(t)$] PV power output at time $t$ (MW)
\item[$\gamma$] Isolation Forest contamination parameter
\end{IEEEdescription}

% =================================================================
\section{Introduction}

Solar irradiance is intrinsically stochastic -- governed by solar geometry, atmospheric attenuation, and cloud movement -- so PV plant output fluctuates unpredictably, jeopardizing voltage stability, power reserves, and dispatch schedules. Accurate ultra-short-term forecasting has therefore become indispensable for grid operators. Deep learning models, particularly LSTM networks \cite{hochreiter1997lstm}, have largely replaced classical statistical methods for solar power prediction due to their ability to capture non-linear temporal dependencies \cite{mellit2021deep, sobri2018review}, and hybrid CNN-LSTM architectures further improve performance by pairing 1D-CNN layers (local ramp dynamics) with LSTM layers (diurnal sequence memory) \cite{zang2020cnnlstm, obiora2021hourly, qu2021dayahead, boucetta2024enhanced}.

Despite these advances, most published solar forecasting literature assumes pristine, noise-free telemetry. In real PV plants this assumption fails: pyranometers, inverters, and SCADA systems regularly suffer \textbf{sensor dropouts} (stuck-at-zero from cable disconnection or inverter shutdown), \textbf{glitch spikes} (multi-fold reading surges from EMI or bit flips), and \textbf{short-duration offset faults} (transient bias from dust accumulation or thermal drift), distorting gradient updates and degrading generalization. While anomaly detection in smart grid telemetry has been studied independently \cite{zhang2021survey}, few studies formalize an end-to-end framework combining unsupervised mitigation with hybrid CNN-LSTM forecasting, isolate detection quality from imputation quality, or subject the resulting claims to formal multiple-comparison and equivalence testing.

\subsection*{Contributions}
\begin{enumerate}
    \item \textbf{End-to-End Fault-Tolerant Architecture:} A unified pipeline coupling unsupervised IF outlier filtering and temporal interpolation with a hybrid 1D CNN-LSTM network.
    \item \textbf{Realistic Fault Injection Protocol:} A reproducible synthetic corruption protocol modeling three physical PV fault modes across $\sim$2\% of telemetry rows.
    \item \textbf{Decomposed Four-Condition Ablation:} Beyond raw-vs-cleaned comparison, we evaluate a \emph{pristine} reference ceiling and an \emph{oracle}-labeled cleaning condition alongside the proposed pipeline, separating detection quality from imputation quality and quantifying what fraction of corruption-induced error is recoverable.
    \item \textbf{Rigorous Multi-Seed Statistical Validation:} All deep learning results are reported as mean $\pm$ standard deviation across 5 random seeds. We report paired $t$-tests for every claimed comparison (including against the pristine and oracle conditions), apply a Holm--Bonferroni correction across the resulting family of tests, and use a TOST equivalence procedure to formally support claims of statistical indistinguishability from the pristine ceiling, rather than inferring equivalence from a non-significant difference alone.
\end{enumerate}

% =================================================================
\section{Related Work}

Early solar forecasting relied on statistical formulations such as ARIMA and exponential smoothing \cite{zhang2015baseline}, which struggle with abrupt ramps and zero-bounded non-linearities, and ML regressors (SVR, Random Forests) that capture non-linear interactions \cite{abumohsen2024hybrid} but lack recurrent mechanisms for sequential dynamics. LSTM and GRU networks \cite{hochreiter1997lstm} overcome the vanishing-gradient problem of vanilla RNNs, and hybrid CNN-LSTM models \cite{zang2020cnnlstm, obiora2021hourly, qu2021dayahead, boucetta2024enhanced, zhou2024enhancing} extend this by using 1D convolutions to extract local shapes (cloud-shadow transitions, morning ramp-ups) ahead of a sequential LSTM stage, sometimes augmented with attention \cite{qu2021dayahead, zhou2024enhancing} or decomposition preprocessing \cite{boucetta2024enhanced} -- directions this work does not pursue, focusing instead on the upstream data-quality problem.

Anomaly detection in power systems has developed through statistical distance checks, autoencoders, and tree-based ensembles \cite{zhang2021survey}. The Isolation Forest \cite{liu2008isolation} offers strong efficiency on high-dimensional time series without distributional assumptions and has been applied to PV fault detection \cite{jorry2026unsupervised}, but is overwhelmingly treated as a standalone monitoring problem \cite{zhang2021survey}; few works integrate it directly into a deep hybrid forecaster's training pipeline or decompose the benefit against oracle-labeled and never-corrupted reference conditions using equivalence testing rather than non-significant $p$-values.

% =================================================================
\section{Methodology}

The pipeline comprises site selection, cyclical feature engineering, synthetic anomaly injection, parallel construction of four data conditions (pristine, raw-corrupted, IF-cleaned, oracle-cleaned), sliding-window sequencing, and CNN-LSTM modeling evaluated identically across all conditions, followed by a multi-seed ($n=5$) benchmark with paired $t$-tests, Holm--Bonferroni correction, and TOST equivalence testing.

\subsection{Dataset and Site Selection}
Data are drawn from the NREL Solar Power Data for Integration Studies (calendar year 2006, SQLite database, 274 PV stations across Alabama). To prevent contamination from physically disparate locations, a strict filter selects a single UPV facility: latitude $30.45^\circ$N, longitude $-88.25^\circ$W; rated capacity $70.0$~MW; sampling resolution $\Delta t = 5$~min; $N = 105{,}120$ consecutive observations (365 days $\times$ 288 samples/day) with zero missing intervals.
\vspace{5pt}
\subsection{Cyclical Astronomical Feature Engineering}
Because the dataset lacks temperature and GHI sensors, four continuous cyclical features are derived from timestamp $t$:
\begin{align}
\text{hour\_sin}(t) &= \sin\!\left(\frac{2\pi h(t)}{24}\right), \;\;
\text{hour\_cos}(t) = \cos\!\left(\frac{2\pi h(t)}{24}\right) \\
\text{day\_sin}(t) &= \sin\!\left(\frac{2\pi d(t)}{365.25}\right), \;\;
\text{day\_cos}(t) = \cos\!\left(\frac{2\pi d(t)}{365.25}\right)
\end{align}
where $h(t) \in [0, 24)$ is fractional hour-of-day and $d(t) \in [1, 365]$ is day-of-year, eliminating the artificial discontinuity at midnight.
\vspace{5pt}
\subsection{Synthetic Anomaly Injection Protocol}
We corrupt $\sim$2.0\% of observations ($N_{\text{anom}} = 2{,}091$ rows) of the pristine series, with three failure types injected in equal proportion:
\begin{enumerate}
    \item \textbf{Sensor Dropout} ($n_1 = 695$): telemetry forced to $0.0$~MW during daylight ($P(t) > 5.0$~MW), simulating severed cables or inverter trips.
    \item \textbf{Glitch Spikes} ($n_2 = 698$): readings multiplied by $\alpha \sim \mathcal{U}(3.0, 5.0)$, clipped to $1.5\times$ capacity, simulating switching surges and ADC bit-flips.
    \item \textbf{Short-Duration Offset Fault} ($n_3 = 698$): blocks of 4--12 consecutive intervals (20--60~min) receive additive offset $\delta \sim \mathcal{U}(0.15, 0.35)\times$ capacity, representing dust accumulation or thermal drift.
\end{enumerate}
The pristine series is retained unmodified in parallel as the reference-ceiling condition.
\vspace{5pt}
\subsection{Anomaly Detection, Oracle Labeling, and Reconstruction}
An IF is trained on $[P_{\text{raw}}(t), \text{hour\_sin}(t), \text{hour\_cos}(t)]$; the path length $h(x)$ across $T=100$ isolation trees yields an anomaly score
\begin{equation}
s(x, n) = 2^{-\frac{\mathbb{E}(h(x))}{c(n)}}
\end{equation}
where $c(n)$ is the average unsuccessful-search path length of a Binary Search Tree. Observations flagged as outliers ($s(x,n) > \tau$, contamination $\gamma = 0.025$) are masked and reconstructed by local linear interpolation:
\begin{equation}
\hat{P}_{\text{clean}}(t) = P(t_a) + \frac{P(t_b) - P(t_a)}{t_b - t_a}(t - t_a)
\end{equation}
where $t_a, t_b$ are the nearest unflagged neighbors; the result is bounded to $[0, P_{\text{max}}]$.

Against ground-truth injection labels, the detector achieved \textbf{72.21\% recall} and \textbf{57.48\% precision} (2{,}627 points flagged). The moderate precision reflects a deliberately permissive setting ($\gamma=0.025$, above the true 2.0\% rate) favoring recall at the cost of false positives near dawn/dusk ramps, where legitimate rapid power changes resemble injected faults.

To separate \emph{detection quality} from \emph{imputation quality}, an \textbf{oracle-cleaned} condition masks the exact injected indices (known a priori) and reconstructs them with the identical interpolation procedure. Comparing IF-cleaned against oracle isolates residual error attributable to imperfect detection, while comparing both against pristine establishes a ceiling neither can exceed.

\subsection{Sliding Sequence Windowing}
Each condition is independently windowed with $W=24$ historical steps (2~h) and horizon $H=6$ steps (30~min):
\begin{align}
\mathbf{X}_\tau &= \begin{bmatrix} \mathbf{x}_{\tau-W+1} & \dots & \mathbf{x}_\tau \end{bmatrix}^T \in \mathbb{R}^{W \times F}, \;\;
y_\tau = P(\tau + H)
\end{align}
with $F=5$ features $[P(t), \text{hour\_sin}, \text{hour\_cos}, \text{day\_sin}, \text{day\_cos}]$, chronologically split 70/15/15 into training ($\approx$73{,}560), validation ($\approx$15{,}766), and test ($\approx$15{,}766) sequences. Min-Max normalization is fitted per-condition on the training partition to prevent leakage; all metrics are computed on de-normalized (MW) values against pristine ground truth.

\subsection{Proposed Hybrid CNN-LSTM Architecture}
\begin{enumerate}
    \item \textbf{1D-CNN Stage:} 64 filters, kernel $k=3$, stride $s=1$, same padding, with batch normalization and ReLU:
    \begin{equation}
    \mathbf{C}_\tau = \text{ReLU}\big(\text{BatchNorm}(\mathbf{W}_c * \mathbf{X}_\tau + \mathbf{b}_c)\big) \in \mathbb{R}^{W \times 64}
    \end{equation}
    \item \textbf{LSTM Stage:} a 2-layer stacked LSTM with hidden dimension $d_h=64$ processes $\mathbf{C}_\tau$:
    \begin{equation}
    \mathbf{h}_t, \mathbf{c}_t = \text{LSTM}\big(\mathbf{C}_\tau(t), \mathbf{h}_{t-1}, \mathbf{c}_{t-1}\big)
    \end{equation}
    \item \textbf{Dense Head:} the final state $\mathbf{h}_W$ is projected through a 32-unit ReLU layer to a single linear output:
    \begin{equation}
    \hat{y}_\tau = \mathbf{w}_o^T\big(\text{ReLU}(\mathbf{W}_d \mathbf{h}_W + \mathbf{b}_d)\big) + b_o
    \end{equation}
\end{enumerate}
Identical architecture and hyperparameters are used for every data condition and for the plain LSTM baseline (which omits the CNN stage), ensuring observed differences are attributable to input condition and/or CNN presence rather than confounding architectural changes.

\subsection{Training Protocol and Multi-Seed Validation}
Training uses Adam ($\eta = 2\times 10^{-3}$, weight decay $10^{-5}$), batch size $B=256$, and MSE loss:
\begin{equation}
\mathcal{L}_{\text{MSE}} = \frac{1}{B}\sum_{i=1}^{B}\left(y_i - \hat{y}_i\right)^2
\end{equation}
Each deep model is trained for 12 epochs with validation-loss checkpointing, repeated across \textbf{5 seeds} $\{42, 123, 7, 2024, 99\}$ governing weight initialization and minibatch shuffling. The underlying datasets (injection, IF detection, oracle labels) are held fixed across seeds so only training stochasticity varies. Persistence and the autoregressive baseline are deterministic and evaluated once.

\subsection{Statistical Validation Protocol}
\label{sec:stat_protocol}
Three paired comparisons are of primary interest: Proposed vs.\ Plain LSTM, vs.\ Pristine Reference, and vs.\ Oracle-Cleaned, each evaluated with a paired $t$-test across the 5 matched seeds. Because three tests are drawn from the same experiment, a \textbf{Holm--Bonferroni correction} controls the familywise error rate. Since a non-significant difference does not by itself establish equivalence, a \textbf{TOST} procedure formally tests whether the Proposed-vs.-Pristine difference falls within a pre-specified margin of $\pm 0.10$~MW ($\approx 0.14\%$ of plant capacity).

% =================================================================
\section{Experimental Results \& Discussion}

\subsection{Performance Metrics}
Performance is evaluated on the de-normalized test set ($N_{\text{test}} \approx 15{,}766$):
\begin{equation}
\text{RMSE} = \sqrt{\tfrac{1}{N}\textstyle\sum_{i=1}^N (y_i - \hat{y}_i)^2}, \quad
\end{equation}
\begin{equation}
\text{MAE} = \tfrac{1}{N}\textstyle\sum_{i=1}^N |y_i - \hat{y}_i|
\end{equation}
\begin{equation}
R^2 = 1 - \frac{\sum_{i=1}^N (y_i - \hat{y}_i)^2}{\sum_{i=1}^N (y_i - \bar{y})^2}
\end{equation}
We additionally define \textbf{Signal Recovery}, the fraction of the corruption-induced error gap a condition closes relative to the pristine ceiling:
\begin{equation}
\text{Signal Recovery} = \frac{\text{RMSE}_{\text{raw}} - \text{RMSE}_{\text{cond}}}{\text{RMSE}_{\text{raw}} - \text{RMSE}_{\text{pristine}}} \times 100\%
\label{eq:sigrecov}
\end{equation}
This is $0\%$ for raw and $100\%$ for pristine by construction. Because the raw-to-pristine denominator is itself small ($\approx$0.31~MW), small absolute RMSE fluctuations produce disproportionate swings in the reported percentage; we treat Signal Recovery as an illustrative point estimate and rely on the paired $t$-tests and TOST procedure of Section~\ref{sec:stat_protocol} -- not the percentage itself -- for equivalence claims.

\subsection{Benchmark Results}
Table~\ref{tab:results} summarizes multi-seed performance across all seven evaluated conditions.

\begin{table*}[!t]
\centering
\caption{Multi-Seed ($n=5$) Performance Benchmark on NREL 70~MW Test Set ($W{=}24$, $H{=}6$, 30-Minute-Ahead Forecast)}
\vspace{-4pt}
\footnotesize{\textit{Deep learning rows report mean $\pm$ standard deviation across 5 seeds $\{42, 123, 7, 2024, 99\}$; Persistence and AutoReg are deterministic and evaluated once. $^\dagger$Signal Recovery is a point estimate (Eq.~\ref{eq:sigrecov}); see Section~IV-A and IV-C2 for the formal statistical tests. $^*$Supported by a TOST equivalence test at a $\pm0.10$~MW margin ($p \approx 0.05$; Section~IV-C2), not by the percentage alone. Proposed-vs.-Plain-LSTM raw paired $t=-3.89$, $p=0.018$; Holm-adjusted $p=0.053$ (Section~IV-C2).}}
\vspace{5pt}
\label{tab:results}
\begin{tabular}{@{}lccccc@{}}
\toprule
\textbf{Model / Data Condition} & \textbf{RMSE (MW) $\downarrow$} & \textbf{MAE (MW) $\downarrow$} & \textbf{$R^2$ $\uparrow$} & \textbf{\% RMSE Impr. vs. Raw} & \textbf{Signal Recovery$^\dagger$} \\
\midrule
CNN-LSTM (Pristine Reference, ceiling) & $3.58 \pm 0.06$ & $1.67 \pm 0.11$ & $0.949 \pm 0.002$ & +7.9\% & 100.0\% (ref.) \\
CNN-LSTM (Oracle-Cleaned)              & $3.58 \pm 0.05$ & $1.67 \pm 0.11$ & $0.949 \pm 0.002$ & +7.9\% & 100.0\%$^*$ \\
\textbf{CNN-LSTM (Proposed, IF-Cleaned)} & $\mathbf{3.57 \pm 0.04}$ & $\mathbf{1.64 \pm 0.08}$ & $\mathbf{0.949 \pm 0.001}$ & \textbf{+8.2\%} & \textbf{100.0\%}$^*$ \\
Plain LSTM (IF-Cleaned)                & $3.77 \pm 0.11$ & $1.78 \pm 0.12$ & $0.943 \pm 0.003$ & +3.0\% & 38.3\% \\
AutoReg (diurnal lag, deterministic)   & 4.83 & 2.76 & 0.907 & $-$24.3\% & 0.0\% \\
Persistence (naive, deterministic)     & 4.93 & 2.35 & 0.903 & $-$26.7\% & 0.0\% \\
CNN-LSTM (Raw Corrupted, ablation)     & $3.89 \pm 0.11$ & $1.97 \pm 0.17$ & $0.940 \pm 0.003$ & Baseline (0.0\%) & 0.0\% (ref.) \\
\bottomrule
\end{tabular}
\end{table*}

\begin{figure*}[!t]
\centering
\includegraphics[width=0.92\textwidth]{figure1_predicted_vs_actual.png}
\caption{Forecast vs.\ ground truth over a representative 24-hour test segment (288 steps, seed 42). The pristine, oracle, and proposed trajectories closely overlap, while the raw-corrupted ablation deviates visibly during ramp regions.}
\label{fig:pred_vs_actual}
\end{figure*}

\FloatBarrier

\subsection{Analysis and Discussion}

\subsubsection{Impact of Anomaly-Aware Preprocessing}
The four-condition ablation isolates preprocessing from architecture; all conditions share identical parameters, hyperparameters, and window structures. The raw-corrupted CNN-LSTM achieves RMSE $=3.89 \pm 0.11$~MW ($R^2 = 0.940$), degraded relative to the pristine ceiling of $3.58 \pm 0.06$~MW ($R^2 = 0.949$). The proposed IF-cleaned model reduces this to $3.57 \pm 0.04$~MW, an 8.2\% mean reduction. This absolute gap is modest because, at $\sim$2\% corruption and a 30-minute horizon, the network already accesses enough uncorrupted context within each 2-hour window to partially compensate for isolated corrupted points; the formal statistical support for how much of this gap is actually closed is given in Section~IV-C2.

Notably, the IF-cleaned and oracle-labeled conditions are nearly identical in mean RMSE ($3.57 \pm 0.04$ vs.\ $3.58 \pm 0.05$~MW), suggesting that at this corruption rate imperfect detection does not materially bottleneck the pipeline relative to knowing exact corruption locations -- the interpolation-based reconstruction itself, rather than detection accuracy, appears to govern recoverable performance, a pattern formally corroborated by the non-significant paired test below.

\subsubsection{Statistical Validation: Multiple Comparisons and Equivalence Testing}
\label{sec:stats_results}

\textbf{1) Paired $t$-tests with Holm--Bonferroni correction.} Table~\ref{tab:stats} reports the three primary paired comparisons across the 5 matched seeds.

\begin{table}[H]
\centering
\caption{Paired $t$-Tests for the Proposed CNN-LSTM, Holm--Bonferroni Corrected}
\label{tab:stats}
\footnotesize
\begin{tabular}{@{}lcccc@{}}
\toprule
\textbf{Comparison} & \textbf{$\Delta$ (MW)} & \textbf{$t$} & \textbf{Raw $p$} & \textbf{Holm $p$} \\
\midrule
vs.\ Plain LSTM   & $-0.2015$ & $-3.892$ & $0.0177$ & $0.0530$ \\
vs.\ Pristine     & $-0.0122$ & $-0.294$ & $0.7837$ & $1.000$ \\
vs.\ Oracle-Cleaned & $-0.0108$ & $-0.280$ & $0.7934$ & $1.000$ \\
\bottomrule
\end{tabular}
\end{table}

The proposed model achieved a lower RMSE than plain LSTM in every one of the 5 paired seeds ($-0.148, -0.204, -0.091, -0.395, -0.170$~MW), and the raw paired test is significant ($p=0.018$). However, correcting for the family of three tests raises the adjusted threshold, and the comparison no longer clears the conventional 0.05 level (Holm-adjusted $p=0.053$) -- a \emph{directionally consistent but not conclusively significant} architectural advantage, more cautious than the uncorrected $p=0.018$ would suggest. Comparisons against Pristine and Oracle-Cleaned remain far from significant after correction (Holm $p=1.000$ for both), consistent with these three conditions performing similarly.

\begin{figure*}[!t]
\centering
\includegraphics[width=0.90\linewidth]{figure2_rmse_bar_chart.png}
\caption{Error metrics across all seven conditions; error bars denote $\pm 1$ standard deviation over 5 seeds. See Table~\ref{tab:stats} for formal significance testing.}
\label{fig:rmse_bar}
\end{figure*}

\textbf{2) TOST equivalence testing.} A non-significant difference does not by itself prove equivalence. We ran a TOST procedure with a pre-specified margin of $\pm 0.10$~MW: mean paired difference $\bar{d} = -0.0122$~MW ($s_d = 0.0928$~MW, $SE = 0.0415$~MW); $t_1 = +2.116$, $p_1 = 0.0509$; $t_2 = -2.703$, $p_2 = 0.0270$; TOST $p = \max(p_1,p_2) = 0.0509$. The 90\% CI for the difference is $[-0.1007, +0.0763]$~MW -- essentially at the conventional 0.05 boundary (at $\pm0.101$~MW, equivalence is confirmed, $p=0.049$). We therefore describe the proposed pipeline as \emph{statistically equivalent to the pristine ceiling within an approximately $\pm 0.10$~MW practical margin}, rather than claiming unqualified 100\% recovery.

\textbf{3) Optimization stability.} The proposed CNN-LSTM shows the lowest RMSE standard deviation ($\pm0.0393$~MW) of any deep model, versus $\pm0.1098$~MW (plain LSTM) and $\pm0.1097$~MW (raw-data ablation), suggesting anomaly-aware cleaning also stabilizes optimization across random initializations.

\subsubsection{Comparison with Classical Baselines}
At 30 minutes ahead, persistence and a diurnal-lag autoregressive model ($t{-}1, t{-}2, t{-}288$) degrade substantially relative to deep approaches (4.93~MW and 4.83~MW, that is $\sim$27\% and $\sim$26\% higher RMSE gaps respectively vs.\ the proposed model), a materially larger margin than at a single-step horizon: at 30 minutes the assumption that the near future resembles the recent past breaks down, and models learning non-linear ramp/ceiling dynamics provide a clear, practically meaningful advantage.

\subsection{Visual Evaluation}
Figs.~\ref{fig:pred_vs_actual}--\ref{fig:anomaly_demo} illustrate forecast-tracking capability relative to pristine and oracle conditions, the aggregate error hierarchy with multi-seed uncertainty, the residual error distribution relative to the ceiling, and a representative anomaly-detection-and-reconstruction example.

\begin{figure}[H]
\centering
\includegraphics[width=0.90\linewidth]{figure3_error_distribution.png}
\caption{Residual error distribution ($\hat{y}-y$, seed 42). The proposed model tracks the pristine ceiling closely; the raw-data ablation shows wider, heavier-tailed dispersion.}
\label{fig:error_dist}
\end{figure}

\begin{figure*}[!t]
\centering
\includegraphics[width=0.88\textwidth]{figure4_anomaly_detection_demo.png}
\caption{Representative 24-hour window where the Isolation Forest detects dropouts and glitch spikes (red `X's) and interpolation reconstructs the diurnal curve. Detector performance over the full test set is 72.21\% recall, 57.48\% precision (Section~III-D); this is a successful case, not a guarantee of flawless detection.}
\label{fig:anomaly_demo}
\end{figure*}

\FloatBarrier

% =================================================================
\section{Conclusion \& Future Work}

We presented an anomaly-aware short-term solar power forecasting framework combining unsupervised IF outlier cleaning with a hybrid 1D CNN-LSTM network at a 30-minute-ahead horizon, evaluated with paired $t$-tests, Holm--Bonferroni correction, and TOST equivalence testing. Decomposing evaluation into pristine, raw-corrupted, IF-cleaned, and oracle-cleaned conditions, the proposed pipeline reaches a Signal Recovery point estimate of 100.0\%, formally supported by a TOST equivalence test showing it is statistically indistinguishable from the pristine ceiling within an approximately $\pm 0.10$~MW margin ($p \approx 0.05$); detection quality does not measurably bottleneck this recovery relative to oracle knowledge (Holm-adjusted $p=1.000$). Across 5 seeds the proposed CNN-LSTM achieves RMSE $=3.57\pm0.04$~MW, outperforming plain LSTM in every paired seed, though after Holm correction this architectural advantage does not reach conventional significance ($p=0.053$) and is reported as consistent but not conclusively confirmed. Both deep models substantially and significantly outperform persistence and autoregressive baselines. At this horizon, preprocessing quality is the more strongly evidenced contributor to accuracy.

\subsection*{Limitations and Future Extensions}
\begin{enumerate}
    \item \textbf{Statistical Power:} $n=5$ seeds underpowers conclusive resolution of the borderline Plain-LSTM comparison ($p=0.053$) and tightening the TOST margin below $\pm0.10$~MW; $n=20$--$30$ is needed.
    \item \textbf{Multi-Site Generalization:} Evaluation uses one 70~MW UPV facility; extension across climates and DPV fleets is needed.
    \item \textbf{Longer Horizons:} Extending from $H=6$ to day-ahead horizons ($H=288$) to test whether the CNN-LSTM's advantage persists.
    \item \textbf{Advanced Imputation:} Bidirectional LSTMs or diffusion-based alternatives to linear interpolation for prolonged multi-hour outages.
\end{enumerate}

% =================================================================
\begin{thebibliography}{00}

\bibitem{hochreiter1997lstm} S. Hochreiter and J. Schmidhuber, ``Long short-term memory,'' \textit{Neural Computation}, vol. 9, no. 8, pp. 1735--1780, 1997.

\bibitem{mellit2021deep} A. Mellit, A. M. Pavan, and V. Lughi, ``Deep learning neural networks for short-term photovoltaic power forecasting,'' \textit{Renewable Energy}, vol. 172, pp. 276--288, 2021.

\bibitem{sobri2018review} S. Sobri, S. Koohi-Kamali, and N. A. Rahim, ``Solar photovoltaic generation forecasting methods: A review,'' \textit{Energy Conversion and Management}, vol. 156, pp. 459--497, 2018.

\bibitem{zang2020cnnlstm} H. Zang et al., ``Short-term solar power prediction based on CNN-LSTM hybrid neural network,'' \textit{Energy Reports}, vol. 6, pp. 217--222, 2020.

\bibitem{obiora2021hourly} C. N. Obiora, A. Ali, and A. N. Hasan, ``Hourly photovoltaic power forecasting using CNN-LSTM hybrid model,'' in \textit{Proc. IEEE PES/IAS PowerAfrica Conf.}, 2021.

\bibitem{qu2021dayahead} J. Qu, Z. Qian, and Y. Pei, ``Day-ahead hourly photovoltaic power forecasting using attention-based CNN-LSTM neural network embedded with multiple relevant and target variables prediction pattern,'' \textit{Energy}, vol. 232, p. 120996, 2021.

\bibitem{boucetta2024enhanced} L. N. Boucetta, Y. Amrane, A. Chouder, S. Arezki, and S. Kichou, ``Enhanced forecasting accuracy of a grid-connected photovoltaic power plant: A novel approach using hybrid variational mode decomposition and a CNN-LSTM model,'' \textit{Energies}, vol. 17, no. 7, p. 1781, 2024.

\bibitem{zhang2021survey} J. E. Zhang et al., ``Time series anomaly detection for smart grids: A survey,'' in \textit{Proc. IEEE Power \& Energy Society Innovative Smart Grid Technologies Conf. (ISGT)}, 2021.

\bibitem{zhang2015baseline} J. Zhang et al., ``Baseline error analysis for solar irradiance and power forecasting,'' \textit{Solar Energy}, vol. 122, pp. 109--124, 2015.

\bibitem{abumohsen2024hybrid} M. Abumohsen, A. Y. Owda, M. Owda, and A. Abumihsan, ``Hybrid machine learning model combining CNN-LSTM-RF for time series forecasting of solar power generation,'' \textit{e-Prime}, vol. 9, p. 100636, 2024.

\bibitem{zhou2024enhancing} N. Zhou, B. Shang, M. Xu, L. Peng, and G. Feng, ``Enhancing photovoltaic power prediction using a CNN-LSTM-attention hybrid model with Bayesian hyperparameter optimization,'' \textit{Global Energy Interconnection}, vol. 7, no. 5, pp. 667--681, 2024.

\bibitem{liu2008isolation} F. T. Liu, K. M. Ting, and Z. H. Zhou, ``Isolation forest,'' in \textit{Proc. IEEE ICDM}, pp. 413--422, 2008.

\bibitem{jorry2026unsupervised} V. Jorry et al., ``An unsupervised kernel norm monitoring for fault detection in a time series photovoltaic system,'' \textit{arXiv preprint} arXiv:2607.01541, 2026.

\bibitem{nrel2016} National Renewable Energy Laboratory (NREL), ``Solar Power Data for Integration Studies,'' Technical Report, NREL/TP-5500-65885, 2016.

\end{thebibliography}

\end{document}
